\documentclass[aspectratio=169,10pt]{beamer}
\usepackage[T1]{fontenc}
\usepackage{lmodern}
\usepackage{amsmath,amssymb}
\usepackage{booktabs}

% --- Look of the course template (same as midterm) ---
\definecolor{tblue}{RGB}{0,51,153}
\definecolor{tred}{RGB}{140,20,20}
\definecolor{keybg}{RGB}{232,240,250}
\usecolortheme[named=tblue]{structure}
\setbeamercolor{frametitle}{bg=tblue,fg=white}
\setbeamercolor{block title}{bg=tblue,fg=white}
\setbeamercolor{block body}{bg=keybg,fg=black}
\setbeamercolor{block title alerted}{bg=tred,fg=white}
\setbeamercolor{block body alerted}{bg=keybg,fg=black}
\setbeamercolor{footline}{bg=tblue,fg=white}
\setbeamertemplate{blocks}[rounded][shadow=true]
\setbeamertemplate{navigation symbols}{}
\setbeamertemplate{itemize items}[triangle]
\setbeamertemplate{frametitle}{%
  \nointerlineskip\begin{beamercolorbox}[wd=\paperwidth,ht=2.6ex,dp=1.2ex,leftskip=0.6em]{frametitle}%
  \usebeamerfont{frametitle}\bfseries\insertframetitle\end{beamercolorbox}}
\setbeamerfont{frametitle}{size=\normalsize}
\setbeamertemplate{footline}{%
  \begin{beamercolorbox}[wd=\paperwidth,ht=2ex,dp=0.8ex,leftskip=0.6em,rightskip=0.6em]{footline}%
  \tiny Trustworthy AI | Final Project Discussion\hfill\insertframenumber/\inserttotalframenumber
  \end{beamercolorbox}}
\newcommand{\keymsg}[1]{%
  \begin{beamercolorbox}[wd=\textwidth,sep=0.7em,rounded=true]{block body}\small\textbf{Key message:} #1\end{beamercolorbox}}

\begin{document}

% ---------------------------------------------------------------- 1
\begin{frame}[plain]
\vspace{1.4em}
\begin{beamercolorbox}[wd=0.92\paperwidth,sep=0.8em,center,rounded=true,shadow=true]{frametitle}
{\large\bfseries Knowing When Not to Trust a Weakly Supervised Model}\\[0.3em]
\small\bfseries Uncertainty-aware abstention under weak supervision and distribution shift
\end{beamercolorbox}
\vspace{1.4em}
\centering
Group 3\\[0.4em] Chlo\'e Deux-Pancera \quad|\quad Tom\'as Saura Mart\'{\i}nez \quad|\quad KACEM HADJI\\[0.6em]
{\scriptsize National Yang Ming Chiao Tung University}\\[0.6em]
Project discussion -- November 11, 2026\\[1.2em]
\footnotesize\textbf{Key papers:} Weak-to-Strong Generalization (A02) \;$\cdot$\; Explainability through Uncertainty (U09)
\;$\cdot$\; Uncertainty-Aware Mixed-Variable ML (U03)
\end{frame}

% ---------------------------------------------------------------- 2
\begin{frame}{Trustworthy AI Issue and Motivation}
\begin{block}{Issue: reliability and human oversight}
A model is often trained on labels from a \textbf{weaker supervisor} (a human, a smaller model) and then
used on \textbf{shifted data}. No ground truth exists to check it, so its \textbf{own uncertainty} is the only
signal for when to defer to a human.
\end{block}
\begin{itemize}\small
\item Our three individual proposals hit the same gap from three sides:
  \begin{itemize}\footnotesize
  \item weak labels (A02) \quad $\cdot$ \quad corrupted inputs (U09) \quad $\cdot$ \quad an unseen group of data (U03)
  \end{itemize}
\item \textbf{Pilot} (CIFAR-10, AlexNet supervisor $\to$ DINO ViT-B/8 student):
  \begin{itemize}\footnotesize
  \item the naive student recovers 53\% of the gap and is under-confident (ECE 8.3\%);
    with the confidence loss: 84\% and ECE 2.7\%, \emph{but}
  \item errors \textbf{copied from the supervisor} are \emph{more} confident than the student's other errors (0.71 vs 0.61).
  \end{itemize}
\end{itemize}
\keymsg{high accuracy does not mean the model knows when it is wrong.}
\end{frame}

% ---------------------------------------------------------------- 3
\begin{frame}{Problem and Research Questions}
\begin{block}{Problem}
Can a weakly supervised model tell us when \emph{not} to trust it, on clean and on corrupted inputs?
\end{block}
\vspace{0.3em}
\begin{enumerate}\small
\item \textbf{Calibration.} Is the student calibrated, and does its uncertainty flag the errors it copied
from the supervisor?
\item \textbf{Abstention.} If the student rejects its most uncertain inputs, does accuracy on the rest
recover, under noise and blur? Which signal works best: softmax, MC Dropout or an ensemble?
\end{enumerate}
\end{frame}

% ---------------------------------------------------------------- 4
\begin{frame}{Key Papers: What We Take and How We Extend}
\footnotesize\setlength{\tabcolsep}{4pt}
\begin{tabular}{p{0.17\textwidth}p{0.34\textwidth}p{0.36\textwidth}}
\toprule
\textbf{Paper} & \textbf{Main method} & \textbf{Our use / extension}\\
\midrule
\textbf{A02} Burns et al., ICML 2024 & Strong student trained on weak labels; PGR; auxiliary confidence loss &
Same setting. \textbf{New:} judge \emph{uncertainty and abstention}, not only accuracy\\[0.4em]
\textbf{U09} Explainability through uncertainty & Uncertainty as explanation; reject uncertain cases;
accuracy--coverage curves & Abstention rule and curves, applied to \emph{weakly supervised} students under shift\\[0.4em]
\textbf{U03} Zhang et al., Sci.\ Rep.\ 2022 & GP / LVGP uncertainty to guide decisions with categorical inputs &
Corruption type as a group: leave-one-corruption-out test of the uncertainty\\
\midrule
Support & \multicolumn{2}{p{0.72\textwidth}}{Ovadia et al.\ 2019 (UQ under shift); CIFAR-10-C (Hendrycks \& Dietterich 2019)}\\
\bottomrule
\end{tabular}
\end{frame}

% ---------------------------------------------------------------- 5
\begin{frame}{Planned Data, Methods and Evaluation}
\begin{columns}[T]
\column{0.5\textwidth}
\begin{block}{Data and models}
\footnotesize
\begin{itemize}
\item \textbf{CIFAR-10} for training; \textbf{CIFAR-10-C} (noise, blur, 5 severities) for shift
\item Weak supervisor: AlexNet / small CNN
\item Strong student: frozen DINO ViT features + small head (cheap)
\item Ceiling: same student on true labels
\item Students: trained on weak labels vs on true labels
\item Uncertainty: softmax, MC Dropout, 5-head ensemble
\end{itemize}
\end{block}
\column{0.5\textwidth}
\begin{block}{Evaluation}
\footnotesize
\begin{itemize}
\item Accuracy and PGR
\item ECE and reliability diagrams, per severity
\item Accuracy--coverage curves (abstention)
\item AUROC of uncertainty for spotting errors, and \emph{copied} supervisor errors
\item 3 seeds; leave-one-corruption-out
\end{itemize}
\end{block}
\end{columns}
\vspace{0.4em}
\footnotesize\textbf{Team split:} weak-to-strong pipeline \;|\; corruptions + abstention \;|\; UQ layer (MC Dropout, ensembles, metrics)
\end{frame}

% ---------------------------------------------------------------- 6
\begin{frame}{Questions for Feedback}
\begin{enumerate}
\item \textbf{Scope.} Is weak supervision \emph{plus} distribution shift too much for one project?
Should we keep only two corruptions (noise, blur)?
\item \textbf{Comparison.} Is a frozen-feature student acceptable, or is full fine-tuning expected?
Are softmax, MC Dropout and ensembles enough as baselines?
\item \textbf{Timeline.} Is this scope realistic for the poster on Dec 23?
Our plan: pipeline and pilot by Nov 18, first results by Dec 2, full results by Dec 16.
\end{enumerate}
\vspace{0.8em}
\keymsg{one pipeline, three angles on the same question: can the model tell us when not to trust it?}
\end{frame}

\end{document}
